\section{問題設定}
\label{sec:problem}


\subsection{予測依存構造}
\(p\)変量時系列を\(\bm x_t=(x_{1,t},\ldots,x_{p,t})^\top\)、最大ラグを\(K\)、利用する過去を\(\bm X_t=(\bm x_{t-1},\ldots,\bm x_{t-K})\)とする。本研究は二乗誤差による予測を扱うため、対象を条件付き平均
\(\mu_i(\bm X_t)=\E[x_{i,t}\mid\bm X_t]\)とし、その近似を\(\widehat x_{i,t}=h_i(\bm X_t)\)と書く。

\begin{definition}[条件付き平均に関する非因果性]
\(\bm X_t^{(-j)}\)を、\(\bm X_t\)から系列\(j\)の全ラグを除いた情報とする。
\begin{equation}
 \E[x_{i,t}\mid\bm X_t]
 =
 \E[x_{i,t}\mid\bm X_t^{(-j)}]
 \quad\text{a.s.}
 \label{eq:mean-noncausality}
\end{equation}
が成立するとき、系列\(j\)の過去は系列\(i\)の条件付き平均の予測を改善しない。
\end{definition}

式\eqref{eq:mean-noncausality}が対象とするのは条件付き平均の依存性である。これに対応するモデル上の性質として、他の入力を固定して系列\(j\)の過去を変えても\(h_i\)が変わらないことを非因果性と定める。

\subsection{線形VARと係数表現}
線形VARの予測式は
\begin{equation}
 \widehat{\bm x}_t=\bm b+\sum_{k=1}^{K}A_k\bm x_{t-k}
 \label{eq:linear-var}
\end{equation}
である。この形式では、全\(k\)について\((A_k)_{ij}=0\)なら系列\(j\)の過去を含む項が全て消え、予測はその系列に依存しない。隣接行列は行を予測対象、列を入力とし、\(\widehat A_{ij}=1\)を\(j\to i\)の向きと定める。

入力依存係数では係数生成を通じた依存も遮断する必要がある。
